% ! TEX root = ../m1-19-20.tex

\chapter{Recapitulare parțial}

\section{Model 1}

\indent\indent 1. Studiați natura seriilor:
\begin{enumerate}[(a)]
\item $ \sum \dfrac{2^n \cdot n!}{n^n} $;
\item $ \sum (\sqrt{n + 1} - \sqrt{n}) $;
\end{enumerate}

\vspace{0.5cm}

2. Studiați convergența uniformă a șirului de funcții:
\[
  f_n : [0, 2] \to \RR, \quad f_n(x) = \dfrac{nx^3}{x^2 + n}.
\]

\vspace{0.5cm}

3. Calculați cu o eroare de maximum $ 10^{-3} $ integrala:
\[
  \int_0^{\frac{1}{3}} \dfrac{\ln(1 + x)}{x} dx.
\]

\vspace{0.5cm}

4. Să se dezvolte în serie Taylor în jurul originii funcția $ f(x) = \ln(2 + 3x) $.

Găsiți domeniul de convergență al dezvoltării.

\vspace{0.5cm}

5. (a) Să se aproximeze, folosind polinomul Taylor de gradul 2 în jurul originii,
funcția $ f(x) = \sqrt[3]{x + 1} $.

(b) Să se calculeze, folosind aproximația de mai sus, $ \sqrt[3]{28} $.

(c) Estimați eroare aproximației de mai sus.

\vspace{0.5cm}

(*) Verificați dacă funcția $ f(x, y) = \arctan\dfrac{x}{y} $ este armonică.

%%%%%%%%%%%%%%%%%%%%%%%%%%%%%%%%%%%%%%%%%%%%%%%%%%%%%%%%%%%%%%%%%%%%%%

\section{Model 2}

\indent\indent 1. Studiați convergența seriilor:
\[
  \sum \dfrac{a^n}{2n^2 + 1}, a \in \RR \qquad \sum \dfrac{(-1)^n}{\ln n}.
\]

\vspace{0.5cm}

2. Studiați convergența punctuală și uniformă a șirului de funcții:
\[
  f_n : [0, 1] \to \RR, \quad f_n(x) = \dfrac{nx^n + 1}{nx^n + 2}.
\]

\vspace{0.5cm}

3. Determinați mulțimea de convergență și suma seriei $ \sum \dfrac{(-1)^n}{3n + 1} x^{3n+1} $.

\vspace{0.5cm}

4. Să se calculeze, cu o eroare mai mică decît $ 10^{-3} $, integrala:
\[
  \int_0^{\frac{1}{3}} \dfrac{\arctan x}{x^2} dx.
\]

\vspace{0.5cm}

5. Găsiți extremele funcției $ f : \RR^3 \to \RR, f(x, y, z) = x^2 + 2z^2 - xy + 2xz - yz + x$.

\vspace{0.5cm}

(*) Calculați aproximația pătratică în jurul originii pentru funcția
$ f(x) = xe^{x^2} $.

%%%%%%%%%%%%%%%%%%%%%%%%%%%%%%%%%%%%%%%%%%%%%%%%%%%%%%%%%%%%%%%%%%%%%%

\section{Parțial 2018--2019}

\textbf{Numărul 1}

\indent\indent 1. Stabiliți natura seriilor:
\begin{enumerate}[(a)]
\item $ \ds\sum_{n \geq 1} \dfrac{\sqrt[3]{n^3 + n^2} - n}{n^2} $;
\item $ \ds\sum_{n \geq 2} (-1)^n \dfrac{n}{n^2 - 1} $.
\end{enumerate}

\vspace{0.5cm}

2. Aproximați cu o eroare de maximum $ \varepsilon = 10^{-3} $ valoarea integralei:
\[
  I = \int_0^1 \dfrac{1 - \cos x}{x^2} dx.
\]

\vspace{0.5cm}

3. Fie seria de puteri: $ \ds\sum_{n \geq 1} \dfrac{(-1)^{n-1}}{n} x^{2n} $.
\begin{enumerate}[(a)]
\item Determinați domeniul de convergență.
\item Calculați suma seriei de puteri.
\end{enumerate}

\vspace{0.5cm}

4. Fie funcția $ f : \RR^2 \to \RR, f(x, y) = ye^{2x + y} $.

Scrieți polinoamele Taylor de gradul 1 și 2 asociate funcției, în jurul punctului
$ A\left( \dfrac{-1}{2}, 1 \right) $.

\vspace{0.5cm}

5. Determinați punctele de extrem local ale funcției:
\[
  f : \RR^3 \to \RR, \quad f(x, y, z) = x^3 + 2y^2 + z^2 + yz - 3x + 6y + 2.
\]

\vspace{1cm}

\textbf{Numărul 2}

1. Studiați natura seriilor numerice:
\begin{enumerate}[(a)]
\item $ \ds\sum_{n \geq 1} \dfrac{\sqrt[3]{n^3 + n^2} - n}{n} $;
\item $ \ds\sum_{n \geq 1} (-1)^n \left( \sqrt{n^2 + 1} - n \right) $.
\end{enumerate}

\vspace{0.5cm}

2. Să se aproximeze cu o eroare de maximum $ 10^{-3}$ valoarea integralei:
\[
  I = \int_0^1 \dfrac{1 - e^{-x}}{x} dx.
\]

\vspace{0.5cm}

3. Fie seria de puteri: $ \ds\sum_{n \geq 0} \dfrac{(-1)^n}{3n + 1} x^{3n + 1} $.
\begin{enumerate}[(a)]
\item Găsiți domeniul de convergență.
\item Calculați suma seriei.
\end{enumerate}

\vspace{0.5cm}

4. Fie funcția $ f : \RR^2 \to \RR, f(x, y) = xe^{x + 2y} $.

Scrieți polinoamele Taylor $ T_1 $ și $ T_2 $ în jurul punctului
$ A\left(1, \dfrac{-1}{2}\right) $.

\vspace{0.5cm}

5. Găsiți punctele de extrem local ale funcției:
\[
  f : \RR^3 \to \RR, \quad f(x, y, z) = x^3 - 2x^2 + y^2 + z^2 - 2xy + xz - yz + 3z.
\]

%%% Local Variables:
%%% mode: latex
%%% TeX-master: "../m1-19-20"
%%% End:
